\partsection{E}{Corrections at the manuscript's uses of envelopes}

We use ``local'' only in the sense of C.3, and ``fibrant'' only for a specified model structure. At IX ms.~9 the letter $\psi$ denotes the restricted Yoneda functor $\Hot\to\Hot_0^{\wedge}$, not the right adjoint asserted at V ms.~37. Its existence is not in dispute. Its proposed detection properties are false. This distinction is necessary for reading the same symbol across the chapters.

\begin{claim}{E.1}{Proved}{V ms.~37--44: the permissible replacements}
The formal propositions on V ms.~37--42 remain theorems with the adjunction hypothesis in A.3--A.4. Their specialization to a right adjoint of $\varphi$ is refuted by B.4. For an individual category $K$, a local category weakly equivalent to $K$ exists if and only if $BK$ is a 1-type; in that case $\eta_K:K\to\Pi_1K$ is such a weak equivalence. For arbitrary $K$, a Thomason fibrant replacement exists, with the mapping property of A.5. For a fixed bounded family of sources, D.1 gives a target that represents the required Hom-functor on that family. For a weak-equivalence-invariant ind-represented functor, A.7 gives its descended ind-object in $\Hot$ without any envelope.
\end{claim}
\begin{proof}
Necessity for the individual local replacement follows from C.3; sufficiency follows because $\Pi_1K$ is local and $B\eta_K$ is a weak equivalence exactly for 1-types. The other assertions have the stated proofs and distinct universal properties. In particular a local category and a fibrant category cannot be substituted for one another in the sentence ``\emph{l'enveloppe de Thomason-Kan}'' at V ms.~42. The unit of the 1-truncation reflection is not a $\W_\infty$-map in the general case.
\end{proof}

\begin{claim}{E.2}{Proved}{V ms.~45: the uniform-stage lemma}
Let $K_0\subset K_1\subset\cdots$ be full subcategories with union $K_\infty$. If $X$ has finitely many objects, with no restriction on its number of arrows, the canonical comparison is an isomorphism of categories
\[
 \colim_n\Fun(X,K_n)\longrightarrow\Fun(X,K_\infty),
\]
and therefore a bijection on components. For an arbitrary filtered diagram of categories the same conclusion holds if $X$ has finitely many arrows. Neither conclusion identifies $\pi_0\Fun(X,K_\infty)$ with $[BX,BK_\infty]$ without an additional argument.
\end{claim}
\begin{proof}
For full inclusions, the finitely many images of objects of $X$ lie at one stage. Fullness puts all the images of arrows there, even if $X$ has infinitely many arrows. A natural transformation has one component for each object of $X$, and its components also lie in a common full stage. Every naturality equation then holds at that stage because its inclusion is faithful. This proves both surjectivity and injectivity on all Hom-sets and objects of the displayed functor. A path in a category is finite, so connectedness of two functors in the union is witnessed at a common stage. For arbitrary filtered diagrams, use the finite presentation argument of D.4. The final caution is witnessed by A.6 and by E.4 below; the comparison to derived maps is a different map.
\end{proof}

\begin{claim}{E.3}{Proved}{V ms.~43 and 46: what fibrant replacement and filtered union supply}
An actual sequence $K_n\to K_{n+1}$ admits a strictly commuting map to a sequence of Thomason-fibrant categories $R_n$ with each $K_n\to R_n$ a trivial cofibration. Its induced map of filtered colimits is in $\W_\infty$. Filtered colimits of Thomason-fibrant categories are Thomason-fibrant. In contrast, filtered unions of local categories need not be local, even with finite full stages.
\end{claim}
\begin{proof}
Choose trivial cofibrations $\alpha_n:K_n\to R_n$ with fibrant targets. The lifting problem whose left side is $\alpha_n$, whose right side is $R_{n+1}\to e$, and whose top composite is $K_n\to K_{n+1}\xrightarrow{\alpha_{n+1}}R_{n+1}$ produces $R_n\to R_{n+1}$. Adjacent squares suffice because the indexing sequence is freely generated. Thomason cofibrations are full sieve inclusions: the corrected statement uses pseudo-Dwyer maps rather than the false assertion that Dwyer maps are closed under retracts; see \citet[Proposition~2.4 and the following discussion]{Raptis2010}.

Nerves commute with filtered colimits, since a simplex and its defining composition equations involve finite data. Weak equivalences of simplicial sets are preserved by filtered colimits; one may verify this by a functorial Kan replacement using finite horns, followed by the finite simplicial definitions of homotopy groups. Hence the colimit of the $\alpha_n$ is a weak equivalence. The generating Thomason horn domains and codomains $c\operatorname{Sd}^2\Lambda_i^r\hookrightarrow c\operatorname{Sd}^2\Delta^r$ are finite categories. A lifting problem into a filtered colimit factors through one stage and receives a filler there, proving the fibrancy assertion. The last assertion is the explicit example in E.4. Thus the sentence at V ms.~46 is true with ``fibrant'' but false with ``local''.
\end{proof}

\begin{claim}{E.4}{Refuted}{V ms.~47--49: an ind-isomorphism need not induce a weak equivalence of unions}
There are finite full filtrations $T_0\subset T_1\subset\cdots$, every $T_n$ local and Thomason-fibrant, with fibrant union $T$, and functors $s,\id_T:T\rightrightarrows T$ such that they induce the same morphism of $(T_n)$ in $\Ind(\Cathot)$, but $Bs$ is not a weak equivalence and is not homotopic to the identity. Thus even the additional fibrancy of every stage and both unions does not rescue either the union-isomorphism inference or the uniqueness question of V ms.~47--49.
\end{claim}
\begin{proof}
Let $G_n$ be the symmetric group on $\{0,\ldots,n\}$, with the inclusions extending permutations by fixed points. Define $T_n$ with objects $0,\ldots,n$ and
\[
 \Hom_{T_n}(i,j)=\begin{cases}G_j&i\leq j,\\\varnothing&i>j.\end{cases}
\]
If $a:i\to j$ has label $g\in G_j$ and $b:j\to k$ has label $h\in G_k$, put $b\circ a=h\,\iota_{jk}(g)$. Group associativity and compatibility of the inclusions prove associativity, and the identity label at each object is its identity arrow. The inclusions $T_n\to T_{n+1}$ are full. The functor $r_n:T_n\to\mathsf BG_n$ sends every object to the unique object and sends labels by inclusion into $G_n$. The inclusion $j_n:\mathsf BG_n\to T_n$ selects object $n$. We have $r_nj_n=\id$ and a natural transformation $\id_{T_n}\Rightarrow j_nr_n$ whose components have identity labels. Hence $T_n$ is isomorphic in $\Cathot$ to a groupoid and is local.

We check the stronger fibrancy assertion separately. For a generating Thomason horn $A\hookrightarrow B$ and $F:A\to T_n$, first extend $r_nF:A\to\mathsf BG_n$ to $B$. Such an extension exists because groupoids are fibrant. Assign every new object of $B$ the object $n$ of $T_n$; keep the old objects and arrows exactly as in $F$. Assign to each arrow involving a new object the label given by the extended groupoid functor. The horn inclusion is a full sieve, so there is no arrow from a new object to an old object. Every remaining arrow with an old source and new target is permitted, because $\Hom_{T_n}(i,n)=G_n$. The label composition rule verifies all mixed composites. This is a strict filler, so $T_n$ is fibrant. E.3 makes $T$ fibrant.

Let $G=\bigcup_nG_n$ and define $s:T\to T$ by $i\mapsto i+1$ and by the shift permutation $s(g)$ which fixes $0$ and carries $g$ on $\{0,\ldots,j\}$ to its action on $\{1,\ldots,j+1\}$. The multiplication rule makes this a functor. For fixed $n$, put $m=n+1$. Both the inclusion $T_n\to T_m$ and the shifted functor $T_n\to T_m$ have natural transformations to the corresponding functors constant on objects at $m$, using identity labels. These last two functors are naturally isomorphic by the label $c=(0\,1\,\cdots\,n+1)$, because
$cgc^{-1}=s(g)$ for $g\in G_n$. Thus the inclusion and shift are connected by three natural transformations already in $\Fun(T_n,T_{n+1})$. Their induced ind-morphisms are equal.

The maps $r_n$ commute with the group inclusions. Filtered-colimit preservation of nerves and homotopy groups, as in E.3, gives $BT\simeq BG$: it is connected, its fundamental group is $G$, and its higher homotopy groups vanish. Under this identification $Bs$ induces the group shift, whose image fixes $0$ and omits $(0\,1)$. It is not surjective on the fundamental group, hence is not a weak equivalence or homotopic to the identity. In particular $T$ is not local. If it were local, the reflection in C.3 would identify its $\Cathot$ endomorphisms with the corresponding $\Hot$ endomorphisms, and E.2 applied to the finite group categories would not itself suffice; directly, the functor $\mathsf BG\to T$ required to represent $\id_{BG}$ would send its object to some finite stage $i$ and factor its entire group map through $G_i$, making a fundamental-group isomorphism impossible. The identity class is therefore missing. The code \texttt{e\_full\_filtration.py} checks every composition and all three transformations through stage three; the uniform formulas above prove the assertion at every stage and in the union.
\end{proof}

\begin{claim}{E.5}{Refuted}{V ms.~48: a local target does not force a strict retract}
Even when $A$ is local, $i:A\hookrightarrow B$ is a Thomason cofibration, and a functor $p:B\to A$ satisfies $[p][i]=\id_A$ in $\Cathot$, there need not exist any functor $r:B\to A$ with $ri=\id_A$ as an equality in $\Cat$.
\end{claim}
\begin{proof}
Let $A$ be the nine-object fence on $0,\ldots,8$, with an arrow from each even vertex to each adjacent odd vertex and no other nonidentity arrows. Let $B$ be the twelve-object circular fence on $\Z/12$, with the same parity rule and adjacency modulo twelve. Include the first nine vertices. This is $c\operatorname{Sd}^2$ applied to the inclusion of a two-edge path in the boundary of a triangle, so it is a Thomason cofibration. The category $A$ is categorically contractible: the successive functors $h_k(i)=\max(i,k)$, $0\leq k\leq8$, are order-preserving and adjacent ones are comparable pointwise. They join the identity to the constant functor at $8$. The constant functor $p:B\to A$ therefore satisfies the required equality in $\Cathot$, and $A$ is local by C.3.

Any functor of these fence posets takes adjacent vertices to equal or adjacent vertices, and hence cannot increase the undirected graph distance. A strict retraction would fix $0$ and $8$. In $B$ they are connected by the four-edge path through $11,10,9$, whereas in $A$ their distance is eight. This is impossible. The exhaustive check of all $9^3=729$ possible choices for the three new vertex images is in \texttt{e\_retraction\_and\_pointed.py}. This example concerns the local meaning of ``cat\'egorie de Thomason''; $A$ is not Thomason-fibrant. Indeed its connected $\operatorname{Ex}^2$-nerve has no edge from $0$ to $8$, because such an edge would be a four-edge-fence map. A connected Kan complex has an edge between every two vertices.
\end{proof}

\begin{claim}{E.6}{Open}{The additional fibrancy version of strict retraction}
For a Thomason-fibrant category $A$, a Thomason cofibration $i:A\hookrightarrow B$, and a functor $p:B\to A$ with $[p][i]=\id_A$ in $\Cathot$, this audit does not decide whether some strict retraction $r$ exists in the same $\Cathot$ class as $p$. The finite example in E.5 does not satisfy this extra fibrancy hypothesis.
\end{claim}
There is always a strict retraction in the same $\Hot$ class as $p$ under these fibrancy hypotheses; no cofibrancy of $A$ is needed for this weaker conclusion. To prove it, encode the specified finite zigzag by $A\times I\to A$, where $I$ is the free category on its finite oriented interval. The constant functor $c:A\to\Fun(I,A)$ is a weak equivalence by the finite interval contraction of E.20. Factor endpoint evaluation as a trivial cofibration $j:\Fun(I,A)\to P$ followed by a fibration $(d_0,d_1):P\to A^2$. Then $jc:A\to P$ is a weak equivalence, $d_0jc=d_1jc=\id_A$, and $A$ fibrant implies that both $d_i:P\to A$ are trivial fibrations. The zigzag supplies $h:A\to P$ with $(d_0h,d_1h)=(pi,\id_A)$. Lift $h$ along the cofibration $i$ against $d_0$, over the map $p:B\to A$. The lift $\widetilde h:B\to P$ gives $r=d_1\widetilde h$ with $ri=\id_A$. In $\Hot$, $d_0=d_1=(jc)^{-1}$, so $[r]=[p]$ there. Nothing in this factorization identifies that model path with a finite natural-transformation zigzag on $B$, which is the remaining requirement in E.6. If $A$ is an actual groupoid the stronger statement is elementary: the zigzag $pi\simeq\id_A$ composes to one natural isomorphism, and conjugating $p$ objectwise by that isomorphism on the full image of $i$, with identity choices elsewhere, gives a strict retraction naturally isomorphic to $p$. These are the verified reductions and the precise unresolved step.

\begin{claim}{E.7}{Proved}{V ms.~49--52: categorical homotopy and a canonical mapping object}
Stagewise natural-transformation homotopies do not imply a global one, even with full fibrant filtrations, by E.4. Without fibrancy the fence of C.5 realizes the manuscript's proposed unbounded-length mechanism. Fibrancy also does not make every categorical homotopy a single transformation. A canonical replacement for arbitrary relative-category mapping types is a derived mapping space of the simplicial localization, whose components are the localized Hom-sets; this construction has no right adjoint $\Hot\to\Cathot$ as an input.
\end{claim}
\begin{proof}
The first two assertions have the indicated explicit counterexamples. For the third take the terminal fork $0\to1\leftarrow2$. Every category with a terminal object is Thomason-fibrant: a functor from a full sieve extends by assigning the terminal object to all new objects and using its unique incoming arrows. This applies to every generating horn. The two functors $e\to C$ selecting $0$ and $2$ are connected by two transformations but by none of length one in either direction. The finite category is checked in \texttt{a\_finite\_checks.py}.

For the mapping object, let $(\mathcal M,W)$ be a relative category in a universe in which its simplicial localization is defined. The hammock localization $L^H(\mathcal M,W)$ has simplicial mapping objects $L^H(\mathcal M,W)(a,b)$ and composition, with
\[
 \pi_0 L^H(\mathcal M,W)(a,b)=\Hom_{\mathcal M[W^{-1}]}(a,b).
\]
This is \citet[\S2 and Proposition~3.1]{DwyerKan1980Calculating}. In a model category these objects compute the derived function complexes \citep[Proposition~4.4 and Corollary~4.7]{DwyerKan1980Function}. These theorems specify a different object from an unqualified ordinary union of strict zigzag categories. Agreement on $\pi_0$ does not give their weak equivalence. A comparison with the manuscript's exact all-arrows path model is proved separately in E.20 and following; for a general $W$ it requires an additional calculus or rectification hypothesis.
\end{proof}

\begin{claim}{E.8}{Proved}{VII ms.~108--124: total categories and actual filtered diagrams}
For an actual filtered diagram $D:I\to\Cat$, the canonical map from its Grothendieck construction to $\colim_I D$ is a weak equivalence. Objectwise weak equivalences of actual filtered diagrams induce weak equivalences of colimits. If $P$ has finite CW type, then
\[
 \colim_i[P,BD_i]\longrightarrow[P,B\colim_iD_i]
\]
is a bijection. These conclusions do not follow from a mere isomorphism of the diagrams as ind-objects in $\Cathot$ or $\Hot$, and the last assertion is false for a general nerve of a finite category.
\end{claim}
\begin{proof}
The nerve of the Grothendieck construction is naturally weakly equivalent to $\hocolim_I ND_i$ by \citet[Theorem~26.8 and Corollary~26.9]{ChacholskiScherer2001}. For a filtered category, homotopy colimit and ordinary colimit of simplicial sets are weakly equivalent: finite simplicial spheres and the finite data specifying their homotopies give the same filtered homotopy groups after a Kan replacement commuting with filtered colimits. The nerve commutes with this colimit by finite presentation of each $[n]$. These observations prove the first two assertions. Compactness of finite CW types proves the third, as in F.7. E.4 refutes the proposed implication from ind-isomorphisms even for full fibrant filtrations. F.1 gives the failure of the last assertion using $BC_2$ and finite skeleta of $BC_2$; for an actual categorical diagram, take $D_r=c\operatorname{Sd}^2\sk_rN(\mathsf BC_2)$. These are finite cofibrant posets, their transition maps are actual functors, and their classifying spaces and colimit are weakly equivalent to the corresponding skeleta and $BC_2$. The distinction invalidates the general inference at VII ms.~121 and its summary on ms.~123. It does not by itself refute a particular geometric comparison for which an independent proof is available.
\end{proof}

\begin{claim}{E.9}{Proved}{External Hom in the rectified setting, with its size condition}
Suppose a relative category $(\mathcal M,W)$ presents an $\infty$-category $\mathcal C$, and suppose, for every test category $X$ under consideration, its strict diagram localization presents the coherent diagram category:
\[
 \mathcal M^X[W_X^{-1}]\simeq h\Fun(NX,\mathcal C),
\]
naturally and compatibly with constant diagrams. Put $Z_{a,b}=\Map_{\mathcal C}(a,b)$. Then
\[
 \Hom_{\mathcal M^X[W_X^{-1}]}(a_X,b_X)
 \cong[BX,Z_{a,b}].
\]
The derived mapping object $Z_{a,b}$, not a Thomason--Kan envelope, is the correct representing object on $\Hot$. If $\mathcal M$ is $\U$-small and the manuscript's countable path ind-representation holds on all $\Cat_\U$, then this additional rectification assertion on all those shapes forces each $Z_{a,b}$ to be a 1-type.
\end{claim}
\begin{proof}
A coherent natural transformation between the constant diagrams at $a,b$ is a coherent $NX$-indexed family of points in $Z_{a,b}$. The end formula for the mapping space, or equivalently the limit of the constant space-valued diagram, gives
\[
 \Map_{\Fun(NX,\mathcal C)}(a_X,b_X)
 \simeq\Map_{\Spc}(BX,Z_{a,b}).
\]
Taking components and applying the stated presentation equivalence proves the formula. This also follows directly, in simplicial sets, from the projective bar resolution $N(X/-)$ used in D.1. A standard sufficient framework for diagram rectification is a combinatorial simplicial model category with its projective diagram structure; the precise coherent-nerve comparison is \citet[Proposition~4.2.4.4]{Lurie2009}. Such a category of all $\U$-small models is generally $\U$-large; one forms it in $\V$. This size distinction is material here.

For a $\U$-small $\mathcal M$, its fixed-length path categories and their countable index form a $\U$-small ind-object. If the displayed Hom-functor equals $[B-,Z_{a,b}]$ on all $\U$-small shapes, D.8 forbids any higher homotopy group of $Z_{a,b}$. Thus the formula for large model presentations, or for bounded shapes, cannot be transported to a small presentation with all shapes without checking this condition. This is also why an arbitrary choice of $W$ does not imply diagram rectification.
\end{proof}

\begin{claim}{E.10}{Proved}{VII ms.~128--136: strict arrows, strict diagrams, and path types}
For every category $C$, the canonical map
\[
 B\Fun([1],C)\longrightarrow\Map(B[1],BC)
\]
is a weak equivalence. Its strict fibre at $(a,b)$ need not be a homotopy fibre. For an arbitrary source $Y$, replacing $C$ by a Thomason-fibrant category does not make the corresponding comparison a weak equivalence, and it need not make $\Fun(Y,C)$ fibrant. A corrected mapping-space formula is $\Map_{\mathbf{sSet}}(NY,RNC)$, or the cosimplicial formula of A.5.
\end{claim}
\begin{proof}
Let $e_0:\Fun([1],C)\to C$ evaluate at the source, and let $c:C\to\Fun([1],C)$ select identity arrows. Then $e_0c=\id_C$ and the square $(\id_x,f)$ at an arrow $f:x\to y$ gives $ce_0\Rightarrow\id$. Thus $Be_0$ is a homotopy equivalence. Evaluation $\Map(B[1],BC)\to BC$ is likewise a homotopy equivalence by contraction of the interval. The canonical comparison commutes with these evaluations, proving the assertion by two-out-of-three. Its strict endpoint fibre is the discrete set $\Hom_C(a,b)$, whereas its derived endpoint fibre is the full path space in $BC$. For the terminal fork, the fibre from $0$ to $2$ is empty although that derived path space is contractible and nonempty. This corrects the absolute assertion at VII ms.~135 while preserving the reason to replace strict endpoint fibres. The two general fibrancy assertions are independently refuted in A.6. No claim that products with arbitrary categories preserve Thomason cofibrations is available: multiplying $\varnothing\to e$ by every $Y$ would already make every category cofibrant.

Grothendieck's statement that the path categories ``\emph{n'ont aucune tendance \`a \^etre des cat\'egories de Thomason-Kan}'' (VII ms.~136) has a precise reading. In the local sense a path category modeling $\Omega Y$ cannot be local whenever $\Omega Y$ has a nonzero homotopy group above degree one. For example paths at a point of a category of type $S^3$ model $\Omega S^3$, whose second homotopy group is $\Z$. No weakly equivalent local category exists by C.3. Fibrancy is a separate property and can be imposed for derived computations; it cannot change that obstruction.
\end{proof}

\begin{claim}{E.11}{Proved}{VII ms.~127, 150--156 and IX ms.~13: what the set-valued arguments prove}
An equality or an inverse verified only in $\Hot_0^{\wedge}$ cannot in general be pulled back to an equality or an inverse in $\Hot$. Thus the finite-test arguments for associativity, $T^2=1$, and weak equivalence of the displayed fibre-transition maps do not prove those conclusions. They are valid statements at the presheaf or ind-object level at which they were established. Fully coherent space-valued finite tests recover the conclusions, and height-two tests detect a specified weak equivalence. Equality of specified maps requires the stronger conditions detailed in Part F. For the actual all-arrows path construction, E.20 and following give direct proofs independent of these invalid inferences.
\end{claim}
\begin{proof}
F.2 gives a specified noninvertible map whose finite-test image is the identity. F.6 gives two unequal maps whose images agree, even with an explicit failure of the proposed epimorphism from an ind-object to its telescope. F.7--F.9 give the positive detection assertions and F.13 types the equality applications. In particular, the argument at VII ms.~152 and 156 only produces an isomorphism of restricted presheaves; conservativity at that level is false. The assertion that $T^2$ acts trivially on an ind-object is likewise insufficient by E.4. None of these counterexamples alone says that the manuscript's particular geometric operation is nonassociative, or that its particular reversal fails to square to one.
\end{proof}

\begin{claim}{E.12}{Proved}{IX ms.~9--13: product-preserving passage and the second psi}
The functors $\varphi$, their ind-extensions, restricted Yoneda $R_1^{\mathrm{cat}}$, and filtered-colimit realization in the presheaf category preserve finite products. They preserve the component sets obtained by evaluation at $e$. Therefore composition already specified as a morphism in one of these categories passes forward and gives the same ordinary category on components. The functor $R_1^{\mathrm{cat}}$, called $\psi$ on this page, is neither faithful, full, nor conservative. Its ind-extension is not faithful or full, since these failures already occur on constant ind-objects.
\end{claim}
\begin{proof}
Finite products of functor categories commute with $\pi_0$: two finite zigzags can be padded by identities to the same alternating length and performed coordinatewise. Realization of nerves preserves finite products up to the natural weak equivalence. Products in an ind-category are represented by pairwise products over the product of the two filtered index categories. Restricted Yoneda preserves products pointwise. Filtered colimits of sets commute with finite products, giving the last preservation assertion for presheaves. Evaluation at $e$ is $\pi_0$ at every level, so ordinary composition, its identities, and its associativity survive these passages. The three failures are F.2--F.3. Since the embedding of a category into its ind-category is fully faithful, the same witnesses between constant ind-objects give the claimed failures for the ind-extension. These facts validate the forward diagram at IX ms.~9 and invalidate the hopes of recovering the higher composition data by reversing its arrows.
\end{proof}

\begin{claim}{E.13}{Refuted}{IX ms.~10: nonconstant ind-objects can have ordinary colimits in Hot}
The assertion that an ind-object of $\Hot$ can have a colimit there only when it is essentially constant is false. The finite discrete spaces $D_n=\{0,\ldots,n\}$ have the discrete countable space $\N$ as colimit in $\Hot$, although $(D_n)$ is not essentially constant in $\Ind(\Hot)$.
\end{claim}
\begin{proof}
For every space $Y$, a map from a discrete space is a choice of a component for each point. Thus
\[
 [\N,Y]=(\pi_0Y)^{\N}=\lim_n(\pi_0Y)^{D_n}=\lim_n[D_n,Y],
\]
naturally, which is the colimit universal property. If the ind-object were represented by a constant $Z$, evaluating its presheaf at $e$ would give $\pi_0Z\cong\N$. The inverse ind-isomorphism $Z\to(D_n)$ would be represented at some finite stage, and composing with the supposed isomorphism back to $Z$ would factor $\id_Z$ in $\Hot$ through $D_n$. This is impossible on components. The existence or nonexistence of a colimit in a homotopy category must therefore be checked for the particular diagram; a homotopy colimit is not automatically that ordinary colimit.
\end{proof}

\begin{claim}{E.14}{Refuted}{IX ms.~22--23: the written tests permit a finite nontrivial relative category}
There is a finite pointed relative category $(\mathcal E,W)$, whose localization is not equivalent to a point, satisfying the representability tests for $\Sigma,\Omega$, the weak-equivalence invariance of the constant-diagram external Hom, and the displayed higher loop-mapping formulas, when those tests are interpreted exactly as written before the finite-limit hypothesis of IX ms.~24. A finite pointed category with all finite products, by contrast, is necessarily equivalent to the terminal category.
\end{claim}
\begin{proof}
Take the full category of pointed sets on the singleton $0$ and the two-element pointed set $x$, and let $W$ contain only identities. The category has five arrows. In particular $\End(x)=\{\id_x,z\}$ with $z^2=z\ne\id_x$; all other Hom-sets are singletons. Thus $0$ is a zero object and the localization is the same nonterminal category. For any small $X$,
\[
 \Hom_{\mathcal E^X}(a_X,b_X)=\Hom_{\mathcal E}(a,b)^{\pi_0X},
\]
since naturality forces the component arrow to be constant on every zigzag component. This inverts every $\W_\infty$-map. The coherent mapping spaces for the ordinary category are discrete, and the loop space of every pointed discrete space is a point. Consequently $\Sigma a=0=\Omega a$ gives both stated representabilities and both mapping-space formulas, with their zero basepoints.

For the actual identity-backward path categories, compose the finite forward chain to define $q:\mathrm{Ch}^W_\infty(a,b)\to\Hom_{\mathcal E}(a,b)_{\mathrm{disc}}$. Insert the composite at a fixed forward edge to define $s$. Replace the part before that edge by $a$, using its accumulated composite into the unchanged tail, to define $R$. Accumulated initial and final composites give natural transformations $R\to\id$ and $R\to sq$, and $qs=\id$. The backward arrows are identities, so all duplicated vertices have equal components and every square commutes by associativity. Thus these exact path categories are isomorphic in $\Cathot$ to discrete categories. The multiplication table and all typed associativity equations of $\mathcal E$ are verified in \texttt{e\_retraction\_and\_pointed.py}.

For the final assertion, if a finite category has two distinct arrows $a\to b$ and all finite powers of $b$, then it has at least $2^r$ arrows $a\to b^r$ for each $r$. This contradicts finiteness of its total arrow set. It is therefore a preorder. In a pointed category every Hom-set is nonempty because of the zero map, so every Hom-set is a singleton and all objects are isomorphic. The extra finite-product hypothesis, rather than tests (a)--(c), supplies this obstruction.
\end{proof}

\begin{claim}{E.15}{Proved}{IX ms.~22--25: representability versus derived loops}
In a pointed $\infty$-category $\mathcal C$ having the required finite limits and colimits, define $\Omega b=0\times_b^h0$ and $\Sigma a=0\amalg_a^h0$. There are canonical pointed equivalences
\[
 \Map_{\mathcal C}(a,\Omega b)\simeq
 \Omega\Map_{\mathcal C}(a,b)\simeq
 \Map_{\mathcal C}(\Sigma a,b),
\]
compatible with loop multiplication. They give $\Sigma\dashv\Omega$ on the homotopy category. In a pointed simplicial model category they can be calculated by replacing $b$ by a fibrant object and factoring its diagonal as
\[
 b\xrightarrow{\simeq}P b\xrightarrow{(d_0,d_1)}b\times b,
\]
with the second arrow a fibration, then taking the fibre at the zero pair. The first arrow must be a weak equivalence; merely calling it a cofibration is insufficient.
\end{claim}
\begin{proof}
The universal property of a homotopy pullback identifies maps into $0\times_b^h0$ with the homotopy pullback $*\times_{\Map(a,b)}^h*$, which is the based loop space. The dual pushout gives the second equivalence. The definitions use the same loop concatenation, so the comparison preserves multiplication and all its coherences. Taking components gives the two representabilities and their adjunction by Yoneda; an arbitrary component-level representability would not conversely prove these space-level equivalences.

For the model calculation take a cofibrant source $a$. The simplicial mapping functor $\Map(a,-)$ preserves the fibration and the defining strict pullback, and carries the weak equivalence between fibrant objects to a weak equivalence. It therefore supplies the corresponding path-object calculation of the loop space. If the first arrow is not required to be a weak equivalence, take a nontrivial Kan $b=K(C_2,1)$, set $P b=b\times b$, and take the second arrow to be the identity. The diagonal is a cofibration of simplicial sets and all stated fibrancy conditions hold, but its strict fibre at the zero pair is a point, whereas $\Omega b$ has two components. This is the exact missing hypothesis in the proposed factorization at IX ms.~24--25.
\end{proof}

\begin{claim}{E.16}{Proved}{The scope of the page corrections}
Every use of the envelope or of the finite-test inference in the specified ranges is assigned a replacement in the following table. A reference to a proved conditional theorem retains its displayed hypotheses; it does not assert that the manuscript has already established them.
\end{claim}
\begin{center}\small
\begin{longtable}{L{0.17\linewidth}L{0.43\linewidth}L{0.31\linewidth}}
\toprule
Manuscript location & Assertion or use & Correct statement and proof\\\midrule
\endfirsthead
\toprule Manuscript location&Assertion or use&Correct statement and proof\\\midrule\endhead
V 37 & Right adjoint asserted; editorial request for a proof & False, B.4 and B.11; positive range C.2.\\
V 37--40 & Reflective-localization Proposition 1 and representability corollary & A.3 under its actual adjunction hypothesis; one-way descent A.7 without it.\\
V 40--42 & Ind-adjunction, Proposition 2, units on ordinary objects & A.4; the ind-unit is tested by the extended functor.\\
V 42--44 & Envelopes and compatible transitions & E.1 and E.3; bounded targets D.1.\\
V 44--46 & External Hom ind-object versus ordinary union & A.7, E.2, E.9; the union need not be local.\\
V 46 & Invariant small ind-objects have local-stage presentations & Exact remaining implication and positive criterion E.17.\\
V 46--47 & Local stages and weak equivalence of union maps & E.4 refutes the inference from local fibrant stages; genuinely local unions satisfy E.18.\\
V 48 & Rectification of a categorical retraction & Local version false, E.5; extra-fibrancy version E.6.\\
V 48--49 & Stagewise homotopies and uniqueness of union maps & Fibrancy alone does not suffice, E.4; genuinely local unions satisfy E.18.\\
V 49 & Shortening a homotopy to one arrow in a fibrant target & False, terminal fork in E.7.\\
V 50--52 & Canonical external homotopy type; proposed return through $\psi$ & E.7 and E.9; no global return through an adjoint.\\
V 52 & Finite tests faithful, full, or conservative & All false for set-valued tests, F.2--F.4.\\
VII 108--115 & Total category, telescope, and all-type projection & Ordinary-colimit leg E.8; exact general all-type projection retained in E.26.\\
VII 116--120 & Independence of arbitrary cofinal systems and transitive comparisons & Canonical collar comparisons E.23 and E.25; the fully quantified original comparisons remain E.26.\\
VII 120--124 & Finite categories compact; conservativity; canonical union & General compactness and detection steps false, F.1--F.2 and E.4; correct compact tests F.7.\\
VII 125--127 & Chosen concatenation, arbitrary blocks, and associativity & Canonical all-arrows operation E.24; relative units E.25; remaining general block and relative assertions E.27.\\
VII 128--130 & Pass to an envelope; reduce to direct one-step paths & E.1, E.9--E.10; strict endpoint arrows do not model general path fibres.\\
VII 131--136 & Strict-diagram and mapping-category arguments & Require diagram rectification E.9; A.6 refutes both fibrancy shortcuts.\\
VII 135 & Absolute arrow-category comparison & True for $[1]$, E.10; the fibre inference is false.\\
VII 136 & Path categories as Thomason--Kan categories & Localness means a 1-type, C.3 and E.10; fibrancy is different.\\
VII 150 & Infer $T^2=1$ by faithfulness & Inference false, E.11; the actual translation satisfies E.22 and its relative extension E.25.\\
VII 150--152 & Endpoint transport inferred from restricted presheaves & Inference false, E.11; direct all-arrows lift-category proof E.21.\\
VII 153--156 & Mapping-path transport and slice transition maps & In the all-arrows scope declared at ms.~137, use E.21 and Quillen B; set-valued conservativity is unavailable.\\
IX 9--10 & Product diagram and the new meaning of $\psi$ & E.12; forward passage and components valid, reverse recovery false.\\
IX 10 & Colimit possible only for essentially constant ind-objects & False, E.13.\\
IX 11--13 & Variance from composition; associativity at ind level and at union & Ind-level Yoneda A.4; E.11, the canonical proof E.24, and the general scope E.27.\\
IX 22--23 & Local path categories identified with invariance; loops and suspension & E.9 and E.15; neither interpretation of ``Thomason--Kan'' can be silently substituted.\\
IX 23 & Proposed finite-category obstruction & Tests alone insufficient, E.14; finite products give the actual obstruction.\\
IX 24--25 & Path-object fibre computes loops & Require the weak-equivalence and fibrancy hypotheses in E.15.\\
\bottomrule
\end{longtable}
\end{center}

\subsection{The remaining assertion about local ind-presentations}

\begin{claim}{E.17}{Open}{V ms.~44 and 46: local presentations of invariant ind-objects}
The following sequential assertion remains unresolved in this audit.
Let $K=(K_n)_{n\in\N}$ be a sequence in $\Cathot_\U$, and assume
that its associated presheaf
\[
 H_K(X)=\colim_n\Hom_{\Cathot}(X,K_n)
\]
sends every map in $\W_\infty$ between $\U$-small categories to a
bijection. Must $K$ be isomorphic in $\Ind_\U(\Cathot_\U)$ to a
sequence of local categories, equivalently to a sequence of groupoids?
The corresponding assertion with an arbitrary $\U$-small filtered
indexing category in place of $\N$ also remains unresolved here.
We do not infer either assertion from descent in A.7, from the
representable-object classification in Part C, or from the obstruction
in D.8 for a single represented homotopy type.
\end{claim}

There is, however, a proved exact criterion for both versions. Let
$\mathcal G\subset\Cathot_\U$ be the full subcategory on groupoids,
let $j:\mathcal G\hookrightarrow\Cathot_\U$ be its inclusion, and
let $\eta_X:X\to j\Pi_1X$ be the unit of the reflection in C.4.
For any $\U$-small filtered $I$ and $K:I\to\Cathot_\U$, the
following conditions are equivalent:
\[
\begin{gathered}
 H_K(\eta_X)\text{ is bijective for every }X\in\Cat_\U;\\
 K\longrightarrow(\Ind_\U j)(\Ind_\U\Pi_1)K
       \text{ is an isomorphism};\\
 K\text{ is isomorphic to a }\U\text{-small ind-object of groupoids}.
\end{gathered}
\tag{E.17a}
\]
These conditions are also equivalent to invariance under $\W_1$,
the functors inducing bijections on components and isomorphisms on
fundamental groups. When they hold, the groupoid presentation can
be taken to be the canonical diagram $(\Pi_1K_i)_{i\in I}$ on the
\emph{same} indexing category. In particular the sequential assertion
does not require an unproved change from a general filtered index
back to a sequential one.

\begin{proof}[Proof of the criterion and the verified special cases]
C.4 gives the fully faithful right adjoint $j$ to $\Pi_1$ and
identifies its inverted class with $\W_1$. Apply A.4 to this actual
adjunction. It gives (E.17a), its $\W_1$ formulation, and the
specified unit. A diagram of local categories is equivalent to a
diagram of groupoids by applying its objectwise invertible units,
using C.3. The canonical comparison on presheaves is explicitly
\[
 \colim_i\Hom_{\Cathot}(X,K_i)
 \longrightarrow
 \colim_i\Hom_{\mathcal G}(\Pi_1X,\Pi_1K_i).
 \tag{E.17b}
\]

If the initial hypothesis of E.17 holds and the $BK_i$ are
$1$-types on a cofinal subdiagram, A.7 identifies the left side
of (E.17b) with $\colim_i[BX,BK_i]$. Each $B\eta_{K_i}$ on
that subdiagram is a weak equivalence, and groupoid locality
identifies the right side with the same colimit. Thus (E.17b)
is bijective and the conclusion holds in this case. Essentially
constant ind-objects with invariant presheaves are also covered
directly by A.7 and C.3, after replacing the diagram by a constant
presentation. More generally, if the descended presheaf on $\Hot_\U$ is
representable by a space $Y$, D.8 forces $Y$ to be a $1$-type.
Then $H_K$ is represented by its fundamental groupoid, so the
ind-object is itself essentially constant with a local representative.
\end{proof}

The missing implication is therefore precisely
\[
 \boxed{\quad
 H_K\text{ inverts }\W_\infty
 \quad\Longrightarrow\quad
 H_K\text{ inverts every }\eta_X:X\to\Pi_1X.
 \quad}
 \tag{E.17c}
\]
The reverse implication follows from (E.17a) because
$\W_\infty\subset\W_1$. The maps $\eta_X$ in (E.17c) are
not weak homotopy equivalences when $BX$ has higher homotopy.
Consequently A.4 applied to the genuine groupoid reflection does
not prove (E.17c) from its stated premise. A.7 only proves
$H_K'\cong\colim_i[-,BK_i]$ after descent; the target varies
with $i$, so D.8 cannot be applied to one of the stages as if it
represented that whole descended presheaf.

The size restriction is material. If the index is allowed to be
$\V$-small and unbounded in the $\U$-cardinals, D.7 gives an
invariant ind-presentation of $X\mapsto[BX,K(\Z/2,2)]$.
It cannot be isomorphic to any ind-object of groupoids: every
such ind-object sends the projection from a finite poset of type
$S^2$ to $e$ to a bijection, whereas this functor takes that map
to the nonbijective constant-map function
$\{*\}\to H^2(S^2;\Z/2)=\Z/2$. This is a counterexample to
the version with the enlarged index universe, and does not answer
the $\U$-small or sequential assertions in E.17. Likewise E.4
refutes the later inference about unions even for a diagram already
made of local categories; it does not refute the existence question
posed in E.17.

\begin{claim}{E.18}{Proved}{V ms.~47--49 with genuinely local unions}
Let $(K_n)$ and $(K'_n)$ be sequences of full inclusions of $\U$-small categories, with unions $K,K'$. Suppose that the unions themselves are local. Let $u:K\to K'$ be an actual functor carrying each $K_n$ into some $K'_{m(n)}$. If the induced morphism in $\Ind_\U(\Cathot_\U)$ is invertible, then $[u]$ is invertible in $\Cathot$. If two such functors $u,v:K\rightrightarrows K'$ induce the same ind-morphism, then $[u]=[v]$ in $\Cathot$. Localness of the unions is a separate hypothesis here; it does not follow from localness of the stages.
\end{claim}
\begin{proof}
For every $\U$-small group $H$, its one-object group category $\mathsf BH$ has finitely many objects, regardless of the size of $H$. E.2 therefore identifies
\[
 \Hom_{\Cathot}(\mathsf BH,K)
       =\colim_n\Hom_{\Cathot}(\mathsf BH,K_n),
\]
and gives the analogous identity for $K'$. These are the actual maps induced by the inclusions. Localness identifies their left sides with $[BH,BK]$ and $[BH,BK']$. An ind-isomorphism induced by $u$ thus gives bijections on these sets for every $H$, and the test $e$ gives a bijection on components.

Both $BK$ and $BK'$ are 1-types by C.3. The preceding bijections respect their decompositions by target component, because restriction along $e\to\mathsf BH$ is evaluation in that component. Every connected 1-type is equivalent to some $BH$. On each corresponding pair of components, ordinary Yoneda on the category of connected 1-types therefore identifies $Bu$ as an isomorphism. These componentwise inverses assemble over the small set of components to an inverse in $\Hot$. By C.3, the full subcategory of local categories maps fully faithfully to the category of 1-types, so that inverse lifts to an inverse of $[u]$ in $\Cathot$.

For the equality assertion, choose an isomorphism in $\Cathot$ from a groupoid of the form $\coprod_{c\in\pi_0K}\mathsf BG_c$ to $K$. Testing each of its component functors with E.2 shows that the two composites with $u,v$ are equal in $\Cathot$ on that component. Compose into the actual groupoid $\Pi_1K'$. A zigzag of transformations between groupoid-valued functors composes, after inverting arrows as needed, to one natural isomorphism. Choose such an isomorphism on every component and assemble them into a single natural isomorphism on the disjoint union. Cancelling the two local groupoidification isomorphisms proves $[u]=[v]$. This explicit assembly deals with infinitely many components and uses no uniform bound on lengths of the original zigzags. The counterexample in E.4 violates the genuine-local-union hypothesis and hence does not contradict this statement.
\end{proof}

\begin{claim}{E.19}{Known}{A coherent replacement for arbitrary relative path composition}
Let $\mathcal M$ be $\U$-small and let $W$ be a subcategory containing every object. The hammock localization, followed on each mapping simplicial set by $\operatorname{Ex}^{\infty}$, gives $\U$-small Kan complexes
\[
 \mathcal H_W(a,b)=\operatorname{Ex}^{\infty}
          L^H(\mathcal M,W)(a,b)
 \qquad(a,b\in\mathcal M)
\]
with strictly associative, strictly unital simplicial composition. Their component category is canonically $\mathcal M[W^{-1}]$. Precomposition or postcomposition by any specified arrow in $W$ is a weak equivalence of these mapping complexes. These are coherent mapping objects for the relative localization. No comparison with a specified ordinary colimit of the manuscript's path categories is part of this assertion.
\end{claim}
\begin{proof}
The hammock simplices, their faces, and their composition are those of \citet[Definition~2.1]{DwyerKan1980Calculating}: a simplex is a finite commuting hammock, its leftward and vertical arrows lie in $W$, and composition concatenates horizontally and reduces consecutive same-direction columns and identity columns. The resulting simplicial category has strictly associative composition and identity hammocks. Every set of finite hammocks lies in $\U$, and taking the countable union of lengths and simplicial degrees preserves that size. Proposition~3.1 of the same paper identifies its component category with the ordinary localization. Proposition~3.3 proves the weak-equivalence assertion for composition with an arrow of $W$.

The functor $\operatorname{Ex}$ preserves products because it is a right adjoint; filtered colimits of sets commute with finite products, so its sequential iteration $\operatorname{Ex}^{\infty}$ also preserves finite products. Applying it to the composition maps therefore gives
\[
 \mathcal H_W(b,c)\times\mathcal H_W(a,b)
       \longrightarrow\mathcal H_W(a,c)
\]
with the same strict associativity and unit equations. The natural maps from a simplicial set to its $\operatorname{Ex}^{\infty}$ replacement are weak equivalences and the replacements are Kan. Thus the component category and the specified transport equivalences are unchanged. In a model-category presentation these are the derived function complexes by \citet[Proposition~4.4 and Corollary~4.7]{DwyerKan1980Function}. This supplies an unconditional relative-category construction in the stated small universe, while comparison with any other raw path model remains a separate mathematical assertion, as recorded in E.26--E.27.
\end{proof}

\subsection{An exact comparison for the unrestricted alternating-path model}

In E.20--E.24 we take $\Sigma=\Mor(X)$. We give the comparison for the precise paths and maps under discussion, rather than replace them silently by another categorical cocylinder. Let $J$ be the integer fence: its objects are $j\in\Z$, its nonidentity arrows are
\[
                 2k\longrightarrow 2k-1,
                 \qquad 2k\longrightarrow 2k+1,
\]
and there are no composable pairs of nonidentity arrows. Put $\mathcal P(X)\subseteq\Fun(J,X)$ for the full subcategory of functors which are constant, with identity arrows, outside some finite interval. A natural transformation between two such functors is necessarily constant on each sufficiently distant tail. Its two eventual components define the endpoint functor
\[
 p=(\sigma,\beta):\mathcal P(X)\longrightarrow X\times X.
\]
Thus $\mathcal P(X)$ is the manuscript's two-sided $\Ch_\infty(X)$ with this explicit parity convention. Its strict fibre $\mathcal P(X;a,b)$ has transformations whose two eventual components are $1_a,1_b$. We write $B$ for classifying space and $K=BX$. All homotopy categories of spaces over a base below mean the localization of the \emph{slice category} at underlying weak equivalences, not the slice of the ordinary homotopy category.

\begin{claim}{E.20}{Proved}{Finite support and the constant-path equivalence}
The constant-path section $j:X\to\mathcal P(X)$ is a weak equivalence. For each $a\in X$, both one-endpoint categories $\mathcal P(X;a,?)$ and $\mathcal P(X;?,a)$ are weakly contractible.
\end{claim}
\begin{proof}
For integers $n\leq m$, let $\mathcal P_{n,m}(X)$ be the full subcategory of paths constant outside $[n,m]$. Restriction identifies it with $\Fun(J|_{[n,m]},X)$. Extension by identities gives fully faithful transition functors when the interval increases. These categories form a filtered union on objects and arrows; every finite diagram is contained in one stage. The same statements hold with either or both endpoints fixed.

A finite oriented interval $I$ with vertices $0,\ldots,r$ has functors $r_k:I\to I$, $r_k(i)=\min(i,k)$, for $0\leq k\leq r$. They preserve each directed generating edge, sending an edge beyond $k$ to an identity. The successive functors $r_k,r_{k+1}$ are joined by a natural transformation in the direction of the edge between $k$ and $k+1$: its components on the tail are that edge, and its other components are identities. On every generating edge the naturality square commutes. Precomposition gives a finite zigzag from the constant-at-initial-value endofunctor of $\Fun(I,X)$ to its identity. This zigzag fixes the initial endpoint. Consequently the constant-diagram functor $X\to\Fun(I,X)$ is a weak equivalence; with initial endpoint fixed at $a$ the functor category is contractible. Reversing the enumeration proves the version with the terminal endpoint fixed.

Nerves preserve filtered colimits of categories, since the data and equations of a composable finite string occur at a stage. Weak equivalences of simplicial sets are preserved by filtered colimits. For completeness, $\operatorname{Ex}^{\infty}$ commutes with filtered colimits: its individual simplex formulas use finite subdivided simplices, and its sequential iteration also commutes with such colimits. In the resulting Kan complexes, homotopy groups and components are detected by finitely many simplices and finite homotopies. Applying this observation to the stagewise constant-diagram equivalences proves the first assertion and the two assertions with one endpoint fixed. No contraction of a strict two-endpoint fibre has been claimed.
\end{proof}

\begin{claim}{E.21}{Proved}{The exact endpoint fibres are homotopy fibres}
For each $a,b\in X$, the natural square
\[
\begin{array}{ccc}
 B\mathcal P(X;a,b)&\longrightarrow&B\mathcal P(X)\\
 \downarrow&&\downarrow Bp\\
 \{(a,b)\}&\longrightarrow&K\times K
\end{array}
\]
is homotopy Cartesian. In particular $B\mathcal P(X;a,b)$ has the weak homotopy type of the ordinary path space $P_{a,b}K$. The comparison with the homotopy fibre is natural for functors over $X\times X$.
\end{claim}
\begin{proof}
We first establish the two lifting properties for this exact category. If $n<m$ are even, both endpoints of $J|_{[n,m]}$ are sources. The endpoint functor on $\Fun(J|_{[n,m]},X)$ is a Grothendieck fibration: to lift $(u,v):(a',b')\to(a,b)$ into a diagram with endpoints $(a,b)$, replace the endpoint objects by $a',b'$ and precompose its two outer arrows by $u,v$. Leave every interior object and arrow unchanged. The transformation to the original diagram has components $u,v$ at the endpoints and identities inside. Any transformation over $(u,v)$ factors uniquely through it, so it is Cartesian. Positive length is essential here. If $n<m$ are odd, both endpoints are sinks, and postcomposition supplies the dual coCartesian construction.

For a general functor $q:E\to B_0$, an object $e\in E$, and an arrow $h:b\to q(e)$, write $L_q(e,h)$ for the category of pairs $(d,\gamma)$ with $q(d)=b$, $q(\gamma)=h$, and $\gamma:d\to e$; its arrows are vertical maps commuting with $\gamma$. For the functor $p$ above, this is the filtered union of the lift categories at sufficiently large positive even-even support intervals. Each has a terminal Cartesian lift, and hence is contractible. Every object, arrow, and equation in a finite diagram of lifts occurs at one such stage. The filtered-colimit argument of E.20 shows that $L_p(e,h)$ is weakly contractible. Odd-odd intervals prove, dually, that every outgoing-lift category is weakly contractible. Both parity families are cofinal among all finite supports. This proves the usual smoothness and properness lift tests directly; it does not assert that the union itself is a Grothendieck fibration or opfibration.

We spell out why these lift tests imply the required Quillen--B condition. Let $q:E\to B_0$ have both properties just proved, and let $E_b$ be its strict fibre. The inclusion
\[
                 i_b:E_b\longrightarrow b\backslash q
\]
is a weak equivalence by Quillen A: its over-comma category at $(e,h:b\to q(e))$ is exactly $L_q(e,h)$. For an arrow $h:b\to b'$, let $M_h$ have as objects arrows $\gamma:d\to d'$ of $E$ over $h$, and as arrows commuting squares whose two vertical sides lie over identities. Let $s:M_h\to E_b$ and $t:M_h\to E_{b'}$ be the source and target projections.

For $d_0\in E_b$, an object of $d_0\backslash s$ is an arrow $\gamma:d\to d'$ over $h$, together with a vertical $\alpha:d_0\to d$. Composition sends it to the outgoing lift $\gamma\alpha:d_0\to d'$. The inclusion of that outgoing-lift category takes $\alpha=1_{d_0}$. The composite in one direction is the identity; in the other, the square $(\alpha,1_{d'})$ gives a natural transformation from the composite to the identity. Thus $d_0\backslash s$ is homotopy equivalent to the contractible outgoing-lift category. Quillen A makes $s$ a weak equivalence. An object of $t/d'_0$ instead consists of $\gamma:d\to d'$ and a vertical $\alpha:d'\to d'_0$. The functor $(\gamma,\alpha)\mapsto\alpha\gamma$ retracts this category onto the incoming lifts into $d'_0$. Its inclusion takes $\alpha=1_{d'_0}$, and the squares $(1_d,\alpha)$ give the natural transformation from the identity to the other composite. Thus $t/d'_0$ is contractible and Quillen A makes $t$ a weak equivalence.

The arrows $\gamma$ themselves give a natural transformation
\[
             i_b s\ \Longrightarrow\ h^*i_{b'}t
             \quad\hbox{into }b\backslash q,
\]
where $h^*:b'\backslash q\to b\backslash q$ is precomposition. As $i_b,s,i_{b'},t$ are weak equivalences, this natural transformation and two-out-of-three show that $h^*$ is a weak equivalence. Quillen B therefore identifies $B(b\backslash q)$ with the homotopy fibre of $Bq$ over $b$, and the natural map from $BE_b$ to that homotopy fibre is a weak equivalence \citep[\S1, Theorems A and B]{Quillen1973}. This formulation is also recorded in \citet[Theorem 1.8]{Shoikhet2016}. The comma constructions and the maps to a standard homotopy-fibre model commute with functors over the base, which gives the asserted naturality.

Apply this argument to $p$. E.20 identifies the constant-path section as a weak equivalence over $K^2$, from the diagonal $K\to K^2$ to $Bp$. Consequently their derived fibres agree. The homotopy fibre of the diagonal over $(a,b)$ is the path space from $a$ to $b$. Equivalently, let $P(K)=K^{[0,1]}$ with its endpoint fibration and constant-path section $s:K\to P(K)$. The comparison is the isomorphism
\[
 u_{\mathcal P}=s\,(Bj)^{-1}:B\mathcal P(X)\longrightarrow P(K)
       \quad\hbox{in }\operatorname{Ho}(\mathrm{Spc}/K^2),
\]
followed by the derived-fibre functor. The preceding proof identifies its source derived fibre with the stated strict categorical fibre.
\end{proof}

\begin{claim}{E.22}{Proved}{The actual even translation is the identity in the homotopy category}
Let $T^2:\mathcal P(X)\to\mathcal P(X)$ be given by $(T^2c)(j)=c(j-2)$, with the corresponding shift on arrows and on natural transformations. It induces the identity in $\Hot$ on the total category and on every fixed-endpoint category.
\end{claim}
\begin{proof}
Translation by two is an automorphism of the integer fence, preserves eventual endpoint values, and fixes every constant diagram and constant natural transformation. Thus $T^2$ is an exact endofunctor over $X^2$ satisfying $T^2j=j$. The morphism $Bj$ is invertible in the localized slice $\operatorname{Ho}(\mathrm{Spc}/K^2)$ by E.20. Cancellation gives $B(T^2)=1$ there. Apply the derived-fibre functor at $(a,b)$. Naturality in E.21 identifies the resulting map with the classifying-space map of the \emph{actual restriction} of $T^2$ to $\mathcal P(X;a,b)$. Hence that map is the identity in $\Hot$. The argument neither assumes faithfulness of finite tests nor constructs a bounded zigzag of natural transformations from the two functors.
\end{proof}

\begin{claim}{E.23}{Proved}{The one-sided models and their actual comparison maps}
Put $D=(+,-)$. Let $\mathcal L(X)$ be the colimit of the categories of $D^n$-paths under padding by a constant $D$-block on the left, and let $\mathcal R(X)$ use padding on the right. They have constant-path sections and endpoint maps. Their constant-path sections are weak equivalences, their one-endpoint fibres are weakly contractible, and their strict two-endpoint fibres compute their derived endpoint fibres. Constant-tail extension gives functors
\[
 q_L:\mathcal L(X)\longrightarrow\mathcal P(X),
 \qquad q_R:\mathcal R(X)\longrightarrow\mathcal P(X)
\]
which are weak equivalences on total categories and on every fixed-endpoint fibre. These comparisons commute with the canonical path-space comparisons of E.21.
\end{claim}
\begin{proof}
Realize $\mathcal L(X)$ on the half-fence $j\leq0$, eventually constant at $-\infty$, using stages $[-2n,0]$; realize $\mathcal R(X)$ on $j\geq0$, using $[0,2n]$. Thus $q_L,q_R$ literally extend the other half by the value at zero. Padding is fully faithful, and all the constant-section and one-endpoint assertions follow from the interval contractions and filtered-colimit proof of E.20.

We prove the comparison on the actual half-fence by sliding an identity collar. Let $S:\mathcal L(X)\to\mathcal L(X)$ append two identity arrows at the right end, shifting the old path two places left. Thus $(Sc)(i)=c(\min(i+2,0))$. We give a homotopy $BS\simeq1$, relative to both endpoints. For $j\geq0$, define an endofunctor $\phi_j$ of the negative half-fence by
\[
 \phi_j(i)=
 \begin{cases}
 i+2,&i\leq-j-2,\\
 -j,&-j-2\leq i\leq-j,\\
 i,&i\geq-j.
 \end{cases}
\]
The values agree at the boundaries. Each generating arrow is either translated by two, unchanged, or sent to an identity; thus $\phi_j$ is a functor and fixes zero. Precomposition gives $F_j(c)=c\phi_j$. We have $F_0=S$. For even $j$ there is a natural transformation $F_j\to F_{j+1}$; for odd $j$ its direction is reversed. Its only possibly nonidentity components are at $-j-2,-j-1$, where both are the edge between $c(-j)$ and $c(-j-1)$. The naturality squares on the three adjacent generating edges commute; all other squares have identity components. The two endpoint components are identities.

For every finite diagram in $\mathcal L(X)$, the functors $F_j$ and these transformations become literally the identity on that diagram for sufficiently large $j$, because its paths and transformations are constant sufficiently far to the left. Concatenate the nerve homotopies on successive time intervals with endpoints $1-2^{-j}$, traversing a natural-transformation homotopy backwards when necessary, and put the identity at time $1$. This is continuous: on each nerve simplex it is stationary after finitely many intervals. In compactly generated spaces, realization has the weak topology with respect to its cells, and product with the compact interval preserves this cellwise test. We obtain $BS\simeq1$, on the total category and on every fixed-endpoint category. A related finite-stage identity-pair sliding calculation appears in \citet[\S3.2.1, Proposition 3.3]{Shoikhet2016}; the explicit formulas above establish the required homotopy for our colimit model independently of a categorical cocylinder identification.

Now $q_LS=T^{-2}q_L$ strictly. Every finite diagram in $\mathcal P(X)$ is carried into the image of $q_L$ by a sufficiently large negative even translation. The same holds for any finite diagram representing a homotopy. Homotopy groups and components of these nerves are computed at finite support by the filtered simplicial argument of E.20. These facts prove that $q_L$ is a weak equivalence, including on every fixed-endpoint fibre. Here is the basepoint detail. For a vertex $c\in\mathcal L(X;a,b)$, the collar homotopy supplies paths from $c$ to $S^kc$ inside that category. Use those paths to identify the relevant based homotopy groups. Translation $T^{-2k}$ gives an automorphism of the homotopy group of $B\mathcal P(X;a,b)$, and the relation $q_LS^k=T^{-2k}q_L$ shows that it preserves the image of the homotopy group of $B\mathcal L(X;a,b)$, in both directions, because $S^k$ is a homotopy equivalence. Any class translates into that image at some finite support, so it was already in the image. If a class from the half-fence becomes zero, translate a finite nullhomotopy into the half-fence; it kills its image under $S^k$, and invertibility of $S^k$ kills the original class. The same argument with paths proves bijectivity on components. Reversal of the fence proves the result for $q_R$.

The total and fibrewise weak equivalences $q_L,q_R$, together with E.21 and naturality of the maps to homotopy fibres, transfer the strict-fibre theorem to $\mathcal L$ and $\mathcal R$. Explicitly, the comparison square has the strict-fibre map on top, the map of homotopy fibres induced by $Bq_L$ or $Bq_R$ below, and the natural fibre comparisons vertically. Its top and bottom maps and its target vertical map are weak equivalences, so its source vertical map is one as well. This argument does not assert that a one-sided endpoint map has both strict categorical lifting properties.

If $j_L,j_R$ denote the constant sections, $q_Lj_L=j$ and $q_Rj_R=j$. Cancellation in the localized slice therefore gives $u_{\mathcal P}Bq_L=s(Bj_L)^{-1}$ and likewise for $q_R$, proving compatibility for these actual maps.

We will also use the elementary iterated-fibre identity. For $(f_1,f_2):A\to B_1\times B_2$, if $A_b=f_1^{-1}(b)\to\operatorname{hofib}_b(f_1)$ is a weak equivalence, regard both spaces over $B_2$ by $f_2$. Their derived fibres over $c$ agree; the latter is $\operatorname{hofib}_{(b,c)}(f_1,f_2)$. This follows from the model consisting of a point of $A$ and a pair of paths in the two base spaces. The maps from strict fibres in this identity are natural.
\end{proof}

\begin{claim}{E.24}{Proved}{The manuscript's left-right concatenation is associative}
For the block $D=(+,-)$ and $\Sigma=\Mor(X)$, glue a left-half path and a right-half path at their common value at zero. This is an actual functor
\[
 \kappa:\mathcal L(X)\times_X\mathcal R(X)\longrightarrow\mathcal P(X),
\]
where the pullback uses the terminal and initial endpoints. On fixed fibres define the manuscript's composition by
\[
 \mu_{a,b,c}=B\kappa\,\bigl((Bq_L)^{-1}\times(Bq_R)^{-1}\bigr):
 B\mathcal P(X;a,b)\times B\mathcal P(X;b,c)
       \longrightarrow B\mathcal P(X;a,c)
 \quad\hbox{in }\Hot.
\]
Here the inverses refer to the fixed-fibre weak equivalences of E.23. Under the canonical comparisons to ordinary path spaces, this \emph{particular} operation is ordinary path concatenation. It is therefore associative and unital in $\Hot$; reversal supplies inverses in the usual path-groupoid sense, and the loop object is a group object of $\Hot$.
\end{claim}
\begin{proof}
The gluing definition agrees on the common vertex and on the common component of a natural transformation, so it defines a functor on the indicated strict pullback. At stages $[-2n,0]$ and $[0,2n]$ it is exactly concatenation of two $D^n$-paths into a $D^{2n}$-path. Increasing $n$ pads at the two outer ends, exactly as in the manuscript's formula. This identifies the specified map, including its transition maps.

Write $E=\mathcal L(X)\times_X\mathcal R(X)$ and let $r:E\to X^3$ record the initial, middle, and terminal values. Its constant-pair section $j_E:X\to E$ is a weak equivalence. One proof uses the filtered system of finite supports: a stage with left length $2n$ and right length $2m$ is the functor category on their glued oriented interval, and its constant section is a weak equivalence by the same explicit interval contraction as in E.20. The constants commute with all enlargements.

We need a relative assertion as well. The middle-coordinate map $BE\to K$ is a weak equivalence since its composite with $Bj_E$ is $1_K$. Its strict fibre at $b$ is
\[
 B\mathcal L(X;?,b)\times B\mathcal R(X;b,?),
\]
which is contractible by E.23. Hence this strict fibre computes the middle-coordinate homotopy fibre. On it the remaining two coordinates are the product of the initial-endpoint map on $\mathcal L(X;?,b)$ and the terminal-endpoint map on $\mathcal R(X;b,?)$. Their strict fibres compute homotopy fibres: apply the iterated-fibre identity in E.23 to each full endpoint map, whose strict-fibre assertion is already proved, using its contractible one-endpoint fibre. Taking the remaining two derived fibres, and using that homotopy fibres commute with products, proves
\[
 B\bigl(\mathcal L(X;a,b)\times\mathcal R(X;b,c)\bigr)
       \simeq \operatorname{hofib}_{(a,b,c)}(BE\xrightarrow{Br}K^3).
\]
Every comparison in this paragraph is the natural comparison from the strict fibre, so it also controls the maps induced by projections and gluing.

Let $P=K^{[0,1]}$ and $Q=P\times_K P$, with their endpoint maps to $K^2$ and $K^3$. Both endpoint maps are fibrations; $Q$ is the space of maps from $[0,2]$ with values specified at $0,1,2$. Write $s:K\to P$ and $t:K\to Q$ for constants. These are weak equivalences. Define in the localized slices
\[
 u_L=s(Bj_L)^{-1},\qquad u_R=s(Bj_R)^{-1},\qquad
 u_{\mathcal P}=s(Bj)^{-1},\qquad v=t(Bj_E)^{-1}.
\]
For each projection from $E$ and for $\kappa$, regard the target over $K^3$ by adjoining the unused endpoint coordinate. This amounts to a product with $K$, which preserves weak equivalences. The two maps comparing categorical gluing with topological concatenation agree after composition with $Bj_E$: each sends a constant pair to a constant path, retaining the same middle coordinate. Since $Bj_E$ is invertible in this slice, cancellation proves
\[
                  u_{\mathcal P}B\kappa
                       =\operatorname{concat}\,v.
\]
The same cancellation proves that $v$ commutes with both projections under $u_L,u_R$.

Apply the derived-fibre functor at $(a,b,c)$. The relative assertion for $E$, E.21, E.23, and the ordinary path fibrations identify the derived fibres with the stated strict fibres. The two projection identities identify the induced map of the product fibre as $u_{L,a,b}\times u_{R,b,c}$, because maps into a product in $\Hot$ are determined by their two projections. Thus
\[
 u_{\mathcal P,a,c}\,B\kappa
 =\operatorname{concat}\,(u_{L,a,b}\times u_{R,b,c}).
\]
Compatibility with $q_L,q_R$ from E.23 gives the claimed formula for $\mu$. Ordinary concatenation is associative up to the endpoint-fixed homotopy which changes the two subdivision points in a threefold concatenation. Constants are units up to the analogous reparameterization homotopy. The reversal $c(j)\mapsto c(-j)$ of the integer fence is an actual functor, swaps endpoints, and preserves the constant section; cancellation in the same slices identifies it with reversal of topological paths. The standard contraction of a path followed by its reverse supplies the inverse identities. These prove all the asserted identities in $\Hot$ for the specified operation.
\end{proof}

\subsection{Relative collars and the remaining quantified comparisons}

We now make explicit the greater generality of VII ms.~108--136. An \emph{admissible} class $\Sigma\subseteq\Mor(X)$ contains every identity and is closed under composition. For a finite sign word $\tau$, let $I_\tau$ be the free category on its oriented interval. A $\Sigma$-path of type $\tau$ is a functor $I_\tau\to X$ taking every geometrically backward arrow to $\Sigma$. The category of such paths is the \emph{full} subcategory of $\Fun(I_\tau,X)$ on those objects. In particular, an interior component of a natural transformation is not required to belong to $\Sigma$; fixed-endpoint transformations have identity components at the two ends. These are the manuscript's definitions at VII ms.~65--67.

\begin{claim}{E.25}{Proved}{Canonical relative collars and literal even translation}
For any admissible $\Sigma$, let $\mathcal L_\Sigma(X;a,b)$, $\mathcal R_\Sigma(X;a,b)$, and $\mathcal P_\Sigma(X;a,b)$ be the full relative subcategories of the canonical $D=(+,-)$ models of E.23. The actual constant-tail extension maps $q_L,q_R$ are weak equivalences on these fibres and on their total categories. The literal translation $(T^2c)(j)=c(j-2)$ induces the identity in $\Hot$ on every such fibre. Gluing with the inverses of $q_L,q_R$ consequently defines the canonical relative operation $\mu_D$ in $\Hot$; constant paths are its left and right units.
\end{claim}
\begin{proof}
Each $\phi_j$ in E.23 is nondecreasing on the geometric vertex order and sends a generating edge to an edge of the same geometric direction or to an identity. Hence $F_j$ preserves relative paths. The homotopies between successive $F_j$ are also permitted, since the relative path category is full and their endpoint components are identities. The eventual-stationarity argument therefore proves $BS\simeq1$ on each relative left-half-fence fibre.

There is an equally explicit right-hand version, without invoking path reversal, which need not preserve a relative path class. For $i,j\geq0$ put
\[
 \psi_j(i)=
 \begin{cases}
 i,&i\leq j,\\
 j,&j\leq i\leq j+2,\\
 i-2,&i\geq j+2.
 \end{cases}
\]
Precomposition gives $G_j$. The functor $G_0=A$ inserts two identities at the left end of a right-half path. Consecutive $G_j$ have their two possibly nonidentity comparison components at $j+1,j+2$, given by the edge between $c(j)$ and $c(j+1)$; the transformation points $G_j\to G_{j+1}$ for even $j$ and backwards for odd $j$. These functors preserve relative paths for the same direction-preservation reason. They become the identity on each finite diagram sufficiently far to the right. Thus $BA\simeq1$, relative to the endpoints.

Even translation preserves the relative class and is an automorphism of the two-sided model. The finite-support and basepoint argument of E.23 now applies word for word, with $q_LS=T^{-2}q_L$ and $q_RA=T^2q_R$, to prove both comparison functors weak equivalences. This use of E.23 requires no endpoint-fibration theorem. Cancelling $Bq_L$ in $Bq_L\,BS=B(T^{-2})Bq_L$ gives $B(T^{-2})=1$, hence $B(T^2)=1$, on the actual relative fibre.

The gluing functor preserves the relative condition, so $\mu_D$ is well defined using the two inverses in $\Hot$. Gluing a left-half path to the constant right-half path is exactly $q_L$; gluing a constant left-half path to a right-half path is exactly $q_R$. Both comparisons preserve constants. Cancellation proves the two unit identities. We have not used the derived-fibre argument of E.24 to claim relative associativity: its strict endpoint-fibre hypothesis was established there only for the all-arrows case.
\end{proof}

To specify the remaining questions, we use the actual type category $\mathrm{Ty}$. Its objects are finite sign words, including the empty word $e$. A morphism is an endpoint-preserving nondecreasing functor between the associated oriented interval categories. Equivalently, if the source word is $\varepsilon_1\cdots\varepsilon_r$, a morphism is a choice of integers $l_i\geq0$ whose target word is $\varepsilon_1^{l_1}\cdots\varepsilon_r^{l_r}$; the $i$th source edge maps to that same-sign string, or to an identity when $l_i=0$. This is the description at VII ms.~50'--53. The \emph{useful length} $\ell_u(\tau)$ is the number of maximal constant-sign runs in the word, with $\ell_u(e)=0$ (VII ms.~75--76).

A projective sequence $t:\N^{\op}\to\mathrm{Ty}$ is \emph{left-cofinal in the manuscript's sense} when
\[
       \forall\tau\in\mathrm{Ty}\quad
       \exists n\quad\Hom_{\mathrm{Ty}}(t_n,\tau)\ne\varnothing.
\]
This is cofinality for the reachability preorder of VII ms.~99--102, not an assumption that comma categories are contractible. Equivalently $\ell_u(t_n)$ is unbounded, as stated at VII ms.~138: a source with sufficiently many alternating sign runs maps to any prescribed word by collapsing unused runs and lengthening the retained same-sign runs; a morphism cannot increase useful length. For such $t$ define the actual category
\[
 C_{\infty,\Sigma}(t;a,b)
   =\colim_n\Fun_\Sigma(I_{t_n},X)_{a,b},
\]
with transition functors given by precomposition with the specified transitions of $t$.

\begin{claim}{E.26}{Open}{The remaining cofinal-type and all-type comparisons}
For arbitrary admissible $\Sigma$ and arbitrary left-cofinal projective type sequences as just defined, this audit leaves undecided the following \emph{specified geometric} assertions of VII ms.~108--123:
\begin{enumerate}
\item A natural transformation $u:t\to t'$ of projective sequences, consisting of strictly commuting maps $u_n:t_n\to t'_n$, always induces a weak equivalence
\[
 u^*:C_{\infty,\Sigma}(t';a,b)\longrightarrow
                  C_{\infty,\Sigma}(t;a,b).
\]
\item The isomorphisms obtained from the canonical type projections $t*t'\to t,t'$ are independent of all comparison choices and form the transitive canonical system asserted at ms.~119--123.
\item The manuscript's comparison from a sequential colimit to its \emph{all-type} path category is an isomorphism in $\Hot$ for every such relative datum. Precisely, let $E_t$ be the fibred Grothendieck construction of the displayed sequence over $\N^{\op}$ and let $\mathrm{Ch}_\Sigma(X;a,b)\to\mathrm{Ty}$ have fibre $\Fun_\Sigma(I_\tau,X)_{a,b}$. The assertion is that
\[
 E_t\longrightarrow\mathrm{Ch}_\Sigma(X;a,b)
\]
is a weak equivalence, so that the span with the ordinary-colimit comparison $E_t\to C_{\infty,\Sigma}(t;a,b)$ yields the claimed isomorphism.
\end{enumerate}
\end{claim}
\begin{proof}[Status and missing hypotheses]
The ordinary-colimit leg of the last span is covered by E.8; its fibred convention is obtained from the covariant one by taking opposite categories. That does not prove its all-type projection is a weak equivalence. The comparison maps in the first two items exist and are the stated precomposition functors, but an isomorphism of their ind-objects does not establish their underlying weak equivalence or equality in $\Hot$; E.4 and E.8 explain the general failure of that inference. E.23 and E.25 establish the concrete collar comparisons, not this entire quantified family of transition systems. No counterexample to these particular geometric assertions, and no complete independent proof of their asserted generality, is supplied here. Thus ``Open'' records the precise remaining audit question, not a claim that the literature has designated it as an open problem.
\end{proof}

\begin{claim}{E.27}{Open}{The remaining relative and arbitrary-block composition assertions}
Fix an admissible $\Sigma$ and a word $\tau_0$ with $\ell_u(\tau_0)\geq2$. Form the three actual systems of VII ms.~125--126: the stages $\tau_0^n$ with transitions which collapse only the extra left blocks, the stages $\tau_0^n$ with transitions which collapse only the extra right blocks, and the stages $\tau_0^{2n}$ with symmetric outer collapse. Denote their colimits by $L_{\tau_0,\Sigma}$, $R_{\tau_0,\Sigma}$, and $P_{\tau_0,\Sigma}$, with endpoints understood. Constant padding defines actual maps $q_L:L_{\tau_0,\Sigma}\to P_{\tau_0,\Sigma}$ and $q_R:R_{\tau_0,\Sigma}\to P_{\tau_0,\Sigma}$, and concatenation defines the actual gluing functor $\kappa_{\tau_0}$. Beyond E.24--E.25, the audit does not settle all of the following claims of ms.~125--127: both padding maps are weak equivalences for every block; the resulting operations $\kappa_{\tau_0}(q_L^{-1}\times q_R^{-1})$ are independent of the block under the specified comparison maps; and those operations are associative in $\Hot$ for every admissible $\Sigma$.
\end{claim}
\begin{proof}[Status and smallest uncovered extensions]
For $\tau_0=D=(+,-)$, E.25 proves that the relative operation exists and is unital. The additional assertion of its associativity for an arbitrary proper admissible class $\Sigma\subsetneq\Mor(X)$ is already an extension not proved here. The proof of E.24 cannot simply be reused: it requires the actual strict relative endpoint fibres to be derived fibres of the diagonal on $BX$, a hypothesis absent for general $\Sigma$. E.9 supplies a different derived mapping-space construction only under its explicit rectification hypothesis and does not identify this raw operation with it.

Even with $\Sigma=\Mor(X)$, the written canonical-block argument does not verify every specified comparison for a different block. An explicit short uncovered convention is the alternating word $(+,-,+)$, whose useful length is three, with the left, right, and symmetric transitions specified in the claim. This is a smallest extension of block length beyond the two-letter canonical case, not a claim of a counterexample. A proof of its comparison and composition compatibility, or a proof for the full quantified family, would settle that extension. The general independence and associativity assertions cannot be inferred from their images under $\psi$. We retain them here with their original quantifiers so that the established all-arrows canonical result does not silently replace them.
\end{proof}

